% GENERATED FILE. Edit canonical lessons, then run npm run generate:books.
% canonical-source-hash: fnv1a64:464a48716b6abc65
% canonical-lessons: MR-C503-visvasghat, MR-C503-bocne, MR-C503-divacne, MR-C503-khadsavne, MR-C503-maph-karne, MR-R503-review-497-503, MR-R503-review-final-503

\chapter{Wounds and Forgiveness}
\label{ch:mr-wounds-and-forgiveness}
\hlchaptermodality{503}

\begin{chapteropening}
\textbf{By the end of this chapter:} \emph{I can talk about a betrayal, words that rankle, a rebuke and forgiving.}

Complete the last lesson of chapter 503: I can talk about a betrayal, words that rankle, a rebuke and forgiving.
\end{chapteropening}

\section[\texorpdfstring{\mr{विश्वासघात} (viśvāsghāt)}{विश्वासघात (viśvāsghāt)}]{\mr{विश्वासघात} (viśvāsghāt) — a betrayal}

\label{lesson:MR-C503-visvasghat}



\begin{quote}
Before the new one: say the Marathi for to seethe, then the Marathi for to stop speaking to.
\end{quote}

\subsection*{You'll want to know: \mr{विश्वासघात}}
\textbf{\mr{विश्वासघात}} — \emph{viśvāsghāt} — \textquotedblleft{}a betrayal\textquotedblright{}.

\begin{quote}
\mr{जवळच्या} \mr{मित्रानेच} \mr{विश्वासघात} \mr{केला}.
\end{quote}

\begin{quote}
\emph{javaḷcyā mitrānec viśvāsghāt kelā.}
\end{quote}

\begin{quote}
It was a close friend who betrayed him.
\end{quote}

Masculine.

\mr{विश्वास}, trust, + \mr{घात}, a blow: \mr{विश्वासघात} \mr{करणे}, to betray.

\begin{cousinweb}
\textbf{\mr{विश्वासघात}} (\emph{viśvāsghāt}) is a tatsama, taken from Sanskrit \emph{viśvāsaghāta}, \textquotedblleft{}destruction of confidence, treachery\textquotedblright{} (Monier-Williams).
\end{cousinweb}

\begin{grammarlens}[title={What you've built}]
One more word for describing what you lived and felt.
\end{grammarlens}

\subsection*{Your turn}
\begin{itemize}
\raggedright
  \item \emph{Say it:} \emph{viśvāsghāt}
  \item \emph{Say it:} \emph{viśvāsghāt}, once more
  \item \emph{Say it:} \emph{javaḷcyā mitrānec viśvāsghāt kelā.} — It was a close friend who betrayed him.
  \item \emph{Recall:} say the Marathi for to seethe, then the Marathi for to stop speaking to, then say \emph{viśvāsghāt} again
\end{itemize}

\begin{tcolorbox}[breakable,colback=teal!4,colframe=teal!35!black,title={Before you move on}]
What does \mr{विश्वासघात} mean? (\textquotedblleft{}A betrayal\textquotedblright{}.) Say it once more.
\end{tcolorbox}
\section[\texorpdfstring{\mr{बोचणे} (bocṇe)}{बोचणे (bocṇe)}]{\mr{बोचणे} (bocṇe) — to rankle; to prick}

\label{lesson:MR-C503-bocne}



\begin{quote}
Before the new one: say the Marathi for to stop speaking to, then the Marathi for a betrayal.
\end{quote}

\subsection*{You'll want to know: \mr{बोचणे}}
\textbf{\mr{बोचणे}} — \emph{bocṇe} — \textquotedblleft{}to rankle; to prick\textquotedblright{}.

\begin{quote}
\mr{त्याचे} \mr{कडू} \mr{शब्द} \mr{मला} \mr{अजून} \mr{बोचतात}.
\end{quote}

\begin{quote}
\emph{tyāce kaḍū śabda malā ajūn boctāt.}
\end{quote}

\begin{quote}
His bitter words still rankle with me.
\end{quote}

Past (to him …): \mr{त्याला} \mr{ते} \mr{बोचलं}, \emph{tyālā te boclaṁ}.

Of thorns and of words; the one hurt takes -\mr{ला}.

\begin{cousinweb}
\textbf{\mr{बोचणे}} (\emph{bocṇe}) is a Marathi verb whose origin is not settled.
\end{cousinweb}

\begin{grammarlens}[title={What you've built}]
One more word for describing what you lived and felt.
\end{grammarlens}

\subsection*{Your turn}
\begin{itemize}
\raggedright
  \item \emph{Say it:} \emph{bocṇe}
  \item \emph{Say it:} \emph{bocṇe}, once more
  \item \emph{Say it:} \emph{tyāce kaḍū śabda malā ajūn boctāt.} — His bitter words still rankle with me.
  \item \emph{Recall:} say the Marathi for to stop speaking to, then the Marathi for a betrayal, then say \emph{bocṇe} again
\end{itemize}

\begin{tcolorbox}[breakable,colback=teal!4,colframe=teal!35!black,title={Before you move on}]
What does \mr{बोचणे} mean? (\textquotedblleft{}To rankle; to prick\textquotedblright{}.) Say it once more.
\end{tcolorbox}
\section[\texorpdfstring{\mr{डिवचणे} (ḍivacṇe)}{डिवचणे (ḍivacṇe)}]{\mr{डिवचणे} (ḍivacṇe) — to goad; to needle}

\label{lesson:MR-C503-divacne}



\begin{quote}
Before the new one: say the Marathi for a betrayal, then the Marathi for to rankle.
\end{quote}

\subsection*{You'll want to know: \mr{डिवचणे}}
\textbf{\mr{डिवचणे}} — \emph{ḍivacṇe} — \textquotedblleft{}to goad; to needle\textquotedblright{}.

\begin{quote}
\mr{भावाने} \mr{मला} \mr{सारखं} \mr{डिवचलं}.
\end{quote}

\begin{quote}
\emph{bhāvāne malā sārkhaṁ ḍivaclaṁ.}
\end{quote}

\begin{quote}
My brother kept needling me.
\end{quote}

Past, with -\mr{ने} (he …): \mr{त्याने} \mr{मला} \mr{डिवचलं}, \emph{tyāne malā ḍivaclaṁ}.

To poke at someone till they react; the one provoked takes -\mr{ला}.

\begin{cousinweb}
\textbf{\mr{डिवचणे}} (\emph{ḍivacṇe}) is a Marathi verb whose origin is not settled.
\end{cousinweb}

\begin{grammarlens}[title={What you've built}]
One more word for describing what you lived and felt.
\end{grammarlens}

\subsection*{Your turn}
\begin{itemize}
\raggedright
  \item \emph{Say it:} \emph{ḍivacṇe}
  \item \emph{Say it:} \emph{ḍivacṇe}, once more
  \item \emph{Say it:} \emph{bhāvāne malā sārkhaṁ ḍivaclaṁ.} — My brother kept needling me.
  \item \emph{Recall:} say the Marathi for a betrayal, then the Marathi for to rankle, then say \emph{ḍivacṇe} again
\end{itemize}

\begin{tcolorbox}[breakable,colback=teal!4,colframe=teal!35!black,title={Before you move on}]
What does \mr{डिवचणे} mean? (\textquotedblleft{}To goad; to needle\textquotedblright{}.) Say it once more.
\end{tcolorbox}
\section[\texorpdfstring{\mr{खडसावणे} (khaḍsāvṇe)}{खडसावणे (khaḍsāvṇe)}]{\mr{खडसावणे} (khaḍsāvṇe) — to rebuke}

\label{lesson:MR-C503-khadsavne}



\begin{quote}
Before the new one: say the Marathi for to rankle, then the Marathi for to goad.
\end{quote}

\subsection*{You'll want to know: \mr{खडसावणे}}
\textbf{\mr{खडसावणे}} — \emph{khaḍsāvṇe} — \textquotedblleft{}to rebuke\textquotedblright{}.

\begin{quote}
\mr{उशिरा} \mr{आल्याबद्दल} \mr{शिक्षकांनी} \mr{त्याला} \mr{खडसावलं}.
\end{quote}

\begin{quote}
\emph{uśirā ālyābaddal śikṣakānnī tyālā khaḍsāvlaṁ.}
\end{quote}

\begin{quote}
The teacher rebuked him for coming late.
\end{quote}

Past, with -\mr{ने} (he …): \mr{त्याने} \mr{मुलाला} \mr{खडसावलं}, \emph{tyāne mulālā khaḍsāvlaṁ}.

Sharper than \mr{रागावणे}; the one rebuked takes -\mr{ला}.

\begin{cousinweb}
\textbf{\mr{खडसावणे}} (\emph{khaḍsāvṇe}) is a Marathi verb whose origin is not settled.
\end{cousinweb}

\begin{grammarlens}[title={What you've built}]
One more word for describing what you lived and felt.
\end{grammarlens}

\subsection*{Your turn}
\begin{itemize}
\raggedright
  \item \emph{Say it:} \emph{khaḍsāvṇe}
  \item \emph{Say it:} \emph{khaḍsāvṇe}, once more
  \item \emph{Say it:} \emph{uśirā ālyābaddal śikṣakānnī tyālā khaḍsāvlaṁ.} — The teacher rebuked him for coming late.
  \item \emph{Recall:} say the Marathi for to rankle, then the Marathi for to goad, then say \emph{khaḍsāvṇe} again
\end{itemize}

\begin{tcolorbox}[breakable,colback=teal!4,colframe=teal!35!black,title={Before you move on}]
What does \mr{खडसावणे} mean? (\textquotedblleft{}To rebuke\textquotedblright{}.) Say it once more.
\end{tcolorbox}
\section[\texorpdfstring{\mr{माफ} \mr{करणे} (māph karṇe)}{माफ करणे (māph karṇe)}]{\mr{माफ} \mr{करणे} (māph karṇe) — to forgive}

\label{lesson:MR-C503-maph-karne}



\begin{quote}
Before the new one: say the Marathi for to goad, then the Marathi for to rebuke.
\end{quote}

\subsection*{You'll want to know: \mr{माफ} \mr{करणे}}
\textbf{\mr{माफ} \mr{करणे}} — \emph{māph karṇe} — \textquotedblleft{}to forgive\textquotedblright{}.

\begin{quote}
\mr{मी} \mr{चूक} \mr{केली}, \mr{तरी} \mr{आईने} \mr{मला} \mr{माफ} \mr{केलं}.
\end{quote}

\begin{quote}
\emph{mī cūk kelī, tarī āīne malā māph kelaṁ.}
\end{quote}

\begin{quote}
I made a mistake, but Mother forgave me.
\end{quote}

Past, with -\mr{ने} (he …): \mr{त्याने} \mr{मला} \mr{माफ} \mr{केलं}, \emph{tyāne malā māph kelaṁ}.

\mr{माफी} \mr{मागणे}, to apologise; the one forgiven takes -\mr{ला}.

\begin{cousinweb}
\textbf{\mr{माफ} \mr{करणे}} (\emph{māph karṇe}) is a Marathi phrase built on \mr{माफ}, \textquotedblleft{}excused\textquotedblright{}, a loan from Persian \emph{muaf}, usually traced to Arabic, with \mr{करणे}.
\end{cousinweb}

\begin{grammarlens}[title={What you've built}]
One more word for describing what you lived and felt.
\end{grammarlens}

\subsection*{Your turn}
\begin{itemize}
\raggedright
  \item \emph{Say it:} \emph{māph karṇe}
  \item \emph{Say it:} \emph{māph karṇe}, once more
  \item \emph{Say it:} \emph{mī cūk kelī, tarī āīne malā māph kelaṁ.} — I made a mistake, but Mother forgave me.
  \item \emph{Recall:} say the Marathi for to goad, then the Marathi for to rebuke, then say \emph{māph karṇe} again
\end{itemize}

\begin{tcolorbox}[breakable,colback=teal!4,colframe=teal!35!black,title={Before you move on}]
What does \mr{माफ} \mr{करणे} mean? (\textquotedblleft{}To forgive\textquotedblright{}.) Say it once more.
\end{tcolorbox}
\section[(dialogue)]{Looking back: family, care and parting}

\label{lesson:MR-R503-review-497-503}



\begin{quote}
Say the Marathi for younger and elder.
\end{quote}

\subsection*{Your turn}
Say these aloud:

\begin{itemize}
\raggedright
  \item younger, elder, the only one, of a paternal uncle, step-
  \item childhood, youth, old age, age, a teenager
  \item pregnant, childbirth, newborn, to crawl, to give to drink
  \item upbringing, to take care of, to nurture, to pamper, to adopt
  \item a touch, to hug, to stroke, to pat, to caress
  \item a separation, bickering, to drift apart, to seethe, to stop speaking to
  \item a betrayal, to rankle, to goad, to rebuke, to forgive
\end{itemize}

\begin{tcolorbox}[breakable,colback=teal!4,colframe=teal!35!black,title={Before you move on}]
Younger? (\textbf{\mr{धाकटा}}, \emph{dhākṭā}.) Elder? (\textbf{\mr{थोरला}}, \emph{thorlā}.) The only one? (\textbf{\mr{एकुलता}}, \emph{ekultā}.) Of a paternal uncle? (\textbf{\mr{चुलत}}, \emph{culat}.) Step-? (\textbf{\mr{सावत्र}}, \emph{sāvatra}.) Childhood? (\textbf{\mr{बालपण}}, \emph{bālpaṇ}.) Youth? (\textbf{\mr{तारुण्य}}, \emph{tāruṇya}.) Old age? (\textbf{\mr{वार्धक्य}}, \emph{vārdhakya}.) Age? (\textbf{\mr{वय}}, \emph{vay}.) A teenager? (\textbf{\mr{किशोर}}, \emph{kiśor}.) Pregnant? (\textbf{\mr{गरोदर}}, \emph{garodar}.) Childbirth? (\textbf{\mr{बाळंतपण}}, \emph{bāḷantpaṇ}.) Newborn? (\textbf{\mr{तान्हा}}, \emph{tānhā}.) To crawl? (\textbf{\mr{रांगणे}}, \emph{rāṅgṇe}.) To give to drink? (\textbf{\mr{पाजणे}}, \emph{pājṇe}.) Upbringing? (\textbf{\mr{संगोपन}}, \emph{saṅgopan}.) To take care of? (\textbf{\mr{सांभाळणे}}, \emph{sāmbhāḷṇe}.) To nurture? (\textbf{\mr{जोपासणे}}, \emph{jopāsṇe}.) To pamper? (\textbf{\mr{लाड} \mr{करणे}}, \emph{lāḍ karṇe}.) To adopt? (\textbf{\mr{दत्तक} \mr{घेणे}}, \emph{dattak gheṇe}.) A touch? (\textbf{\mr{स्पर्श}}, \emph{sparśa}.) To hug? (\textbf{\mr{मिठी} \mr{मारणे}}, \emph{miṭhī mārṇe}.) To stroke? (\textbf{\mr{गोंजारणे}}, \emph{goñjārṇe}.) To pat? (\textbf{\mr{थोपटणे}}, \emph{thopaṭṇe}.) To caress? (\textbf{\mr{कुरवाळणे}}, \emph{kurvāḷṇe}.) A separation? (\textbf{\mr{ताटातूट}}, \emph{tāṭātūṭ}.) Bickering? (\textbf{\mr{कुरबुर}}, \emph{kurbur}.) To drift apart? (\textbf{\mr{दुरावणे}}, \emph{durāvṇe}.) To seethe? (\textbf{\mr{धुसफुसणे}}, \emph{dhusphusṇe}.) To stop speaking to? (\textbf{\mr{अबोला} \mr{धरणे}}, \emph{abolā dharṇe}.) A betrayal? (\textbf{\mr{विश्वासघात}}, \emph{viśvāsghāt}.) To rankle? (\textbf{\mr{बोचणे}}, \emph{bocṇe}.) To goad? (\textbf{\mr{डिवचणे}}, \emph{ḍivacṇe}.) To rebuke? (\textbf{\mr{खडसावणे}}, \emph{khaḍsāvṇe}.) To forgive? (\textbf{\mr{माफ} \mr{करणे}}, \emph{māph karṇe}.)
\end{tcolorbox}
\section[(dialogue)]{Wounds and Forgiveness — once more}

\label{lesson:MR-R503-review-final-503}



\begin{quote}
Say the Marathi for a betrayal and to rankle.
\end{quote}

\subsection*{Your turn}
\begin{itemize}
\raggedright
  \item \emph{Say it:} a betrayal, to rankle, to goad, to rebuke, to forgive
\end{itemize}

\begin{tcolorbox}[breakable,colback=teal!4,colframe=teal!35!black,title={Before you move on}]
A betrayal? (\textbf{\mr{विश्वासघात}}, \emph{viśvāsghāt}.) To rankle? (\textbf{\mr{बोचणे}}, \emph{bocṇe}.) To goad? (\textbf{\mr{डिवचणे}}, \emph{ḍivacṇe}.) To rebuke? (\textbf{\mr{खडसावणे}}, \emph{khaḍsāvṇe}.) To forgive? (\textbf{\mr{माफ} \mr{करणे}}, \emph{māph karṇe}.)
\end{tcolorbox}
